\documentclass[11pt]{amsart}
\usepackage[margin=1.1in]{geometry}
\usepackage{amsmath,amssymb,amsthm,mathtools}
\usepackage{enumitem}
\usepackage{microtype}
\usepackage{hyperref}
\usepackage[nameinlink,capitalize]{cleveref}
\hypersetup{colorlinks=true,linkcolor=blue,citecolor=blue,urlcolor=blue}
\pdfstringdefDisableCommands{%
  \def\({}%
  \def\){}%
  \def\tilde#1{#1}%
  \def\mathrm#1{#1}%
  \def\mathbb#1{#1}%
  \def\tau{tau}%
  \def\lambda{lambda}%
  \def\infty{infty}%
}

\theoremstyle{plain}
\newtheorem{theorem}{Theorem}[section]
\newtheorem{lemma}[theorem]{Lemma}
\newtheorem{proposition}[theorem]{Proposition}
\newtheorem{corollary}[theorem]{Corollary}
\theoremstyle{definition}
\newtheorem{definition}[theorem]{Definition}
\newtheorem{assumption}[theorem]{Assumption}
\theoremstyle{remark}
\newtheorem{remark}[theorem]{Remark}

\title[Assembly of the Local Type II Exclusion]{Assembly of the Local Type II Exclusion Theorem for Navier--Stokes Blowup}
\author{Guillem Cordero}
\date{}

\begin{document}

\begin{abstract}
We assemble the local Type II exclusion argument for the three-dimensional
incompressible Navier--Stokes equations in standard PDE language.  The input is
a suitable weak solution near a putative singular point, the local
Caffarelli--Kohn--Nirenberg concentration theorem, a local Type I tangent
criterion, the local single-core Type II criterion, the localized scale-translation representation,
the local multibubble and cascade exclusion, the compact-cylinder Caccioppoli
regularity theorem, and the scale-collapse dichotomy.  Under the hypotheses
stated in those results, every positive local Type II concentration is exhausted
by the single-core, multibubble/cascade, rough-core, or scale-collapse
alternatives, and each alternative is excluded by the corresponding theorem.
Consequently no Type II singularity occurs at the point under these analytic
hypotheses.
\end{abstract}

\maketitle

\section{Introduction}

Let \((u,p)\) be a suitable weak solution of the three-dimensional
incompressible Navier--Stokes equations
\[
  \partial_tu+(u\cdot\nabla)u+\nabla p=\nu\Delta u,
  \qquad \nabla\cdot u=0,
\]
in a neighborhood of a possible singular point \(z_0=(x_0,T)\).  This paper assembles the local Type II part of the analysis from the local entry theorem and the local Type II components stated below.

The entry point is local.  For \(r>0\), set
\[
  Q_r(z_0)=B_r(x_0)\times(T-r^2,T),
\]
and define the pressure-normalized Caffarelli--Kohn--Nirenberg quantities
\[
  C(z_0,r)=r^{-2}\int_{Q_r(z_0)}|u|^3\,dx\,dt,
\]
\[
  D(z_0,r)=r^{-2}\int_{T-r^2}^T\int_{B_r(x_0)}
  |p(x,t)-(p)_{B_r(x_0)}(t)|^{3/2}\,dx\,dt.
\]
The pressure is always understood modulo functions of time, as in the local
energy inequality and the CKN theory \cite{CaffarelliKohnNirenberg1982}.

The local concentration theorem says that a singular point has positive local
CKN concentration.  After the local Type I tangent criterion removes nonzero
bounded parabolic tangent limits, the remaining positive concentration is the
local Type II regime.  The local Type II analysis divides that regime into the
following PDE alternatives.

\begin{enumerate}[label=\textup{(\roman*)}]
\item A single selected core with bounded compact-cylinder suitable package and
nondegenerate localized scale-translation gauges.
\item A multibubble, cascade, or gauge-degenerate configuration.
\item Failure of local windowed \(H^1\) control in a bounded renormalized core.
\item Collapse of the renormalized scale.
\end{enumerate}

The single-core alternative is excluded by the local compact criterion.  The
multibubble and cascade alternatives are excluded by the local decoupling and
scale-rigidity theorems.  The rough-core alternative is excluded by the
compact-cylinder Caccioppoli theorem.  The scale-collapse alternative is
excluded either by a finite-cost condition controlled by the scale-drift cost or
by reduction to a stationary or generalized self-similar profile in a parameter
regime covered by the relevant Liouville theorem.

\begin{theorem}[Local Type II exclusion, assembled form]\label{thm:intro-main}
Let \((u,p)\) be a suitable weak solution near \((x_0,T)\).  Assume the local
Type I tangent criterion.  Assume also the analytic hypotheses of the five
local PDE components: local single-core exclusion, localized
scale-translation representation, local multibubble/cascade exclusion,
compact-cylinder Caccioppoli regularity, and scale-collapse exclusion by cost or
by stationary-profile rigidity.  Then \((x_0,T)\) is not a Type II singular
point.
\end{theorem}

The theorem is deliberately conditional on the analytic hypotheses named above.
It does not assert global regularity for arbitrary initial data.  Its conclusion
is that, once the local Type I criterion and the five local Type II components
are available in their stated forms, the Type II alternatives are exhausted and
excluded.

\section{Local entry and Type I separation}

\begin{theorem}[Positive local concentration at singular points]\label{thm:positive-concentration}
If \((x_0,T)\) is singular, then
\[
  \limsup_{r\downarrow0}\bigl(C((x_0,T),r)+D((x_0,T),r)\bigr)>0.
\]
Equivalently, there are \(\eta>0\) and \(r_n\downarrow0\) such that
\[
  C((x_0,T),r_n)+D((x_0,T),r_n)\ge \eta
  \qquad\text{for all }n.
\]
\end{theorem}

\begin{proof}
This is the contrapositive of CKN epsilon regularity.  If the displayed limsup
were zero, then for sufficiently small \(r\) the CKN quantity would be below the
universal regularity threshold, and \((x_0,T)\) would be regular.
\end{proof}

Given such a sequence, define the parabolic rescalings
\[
  u_n(y,s)=r_n u(x_0+r_ny,T+r_n^2s),
  \qquad
  p_n(y,s)=r_n^2p(x_0+r_ny,T+r_n^2s).
\]
Then \((u_n,p_n)\) is suitable on every fixed backward cylinder once \(n\) is
large enough, and the unit cylinder retains positive normalized CKN mass.

\begin{assumption}[Local Type I tangent criterion]\label{ass:typeI}
Any blowup sequence whose selected-window rescalings have a nonzero bounded
parabolic tangent on a compact cylinder is excluded by the Type I theory.
\end{assumption}

A positive local concentration sequence not excluded by this criterion is called
a local Type II sequence in this paper.

\section{Compact-cylinder dichotomy}

For a suitable pair \((v,q)\) on a backward cylinder, write
\[
  A_v(R)=R^{-1}\operatorname*{ess\,sup}_{-R^2<s<0}
  \int_{B_R}|v(y,s)|^2\,dy,
\]
\[
  E_v(R)=R^{-1}\int_{-R^2}^0\int_{B_R}|\nabla v|^2\,dy\,ds,
\]
and let \(C_{v,q}(R)\), \(D_{v,q}(R)\) denote the corresponding scaled local
\(L^3\) and pressure quantities.

\begin{theorem}[Local compactness dichotomy]\label{thm:compactness-dichotomy}
After passing to a subsequence of the local rescalings \((u_n,p_n)\), exactly
one of the following alternatives holds.

\begin{enumerate}[label=\textup{(\roman*)}]
\item For every fixed \(R<\infty\),
\[
  \sup_n\left(A_{u_n}(R)+E_{u_n}(R)+C_{u_n,p_n}(R)+D_{u_n,p_n}(R)\right)<\infty.
\]
In this case a subsequence converges locally in the suitable weak class on
compact backward cylinders.
\item For some fixed \(R<\infty\), the preceding supremum is infinite.  Then a
bounded rescaled cylinder carries an additional local concentration, giving a
multibubble or cascade alternative.
\end{enumerate}
\end{theorem}

\begin{proof}
The first alternative is precisely the compact-cylinder suitable package.  If it
holds, the standard compactness theorem for suitable weak solutions gives a
subsequence with local weak suitable convergence, strong convergence in
subcritical local spaces, and pressure convergence modulo spatial means.  If it
fails, then at least one component of the local suitable package is unbounded on
a fixed cylinder.  The local Caccioppoli inequality shows that boundedness of
\(C+D\) on a comparable larger cylinder controls \(A+E\) on a smaller one.
Thus failure of the compact-cylinder package forces additional bounded-cylinder
concentration, which is the local multibubble or cascade alternative.
\end{proof}

\section{Imported PDE components}

This section states the local PDE components used in the assembly.

\begin{theorem}[Local single-core criterion]\label{thm:single-core}
Assume the local Type I tangent criterion.  A local Type II sequence with one
selected core, bounded compact-cylinder suitable package, nondegenerate
localized scale-translation gauges, compact-ball pressure decomposition, local
Caccioppoli bounds, and finite localized Type II cost cannot occur.
\end{theorem}

\begin{proof}
Finite localized cost gives windows on which the localized dissipation tends to
zero.  Compact-cylinder bounds give convergence on the same windows.  The
low-dissipation limit is spatially constant on the selected core.  If the
constant is zero, the selected CKN mass vanishes, contradicting positive local
concentration.  If the constant is nonzero, the sequence has a nonzero bounded
parabolic tangent and is excluded by the local Type I tangent criterion.
\end{proof}

\begin{theorem}[Localized scale-translation representation]\label{thm:representation}
On a genuine single-core local Type II sequence, the localized scale and
centering conditions can be imposed on compact cylinders whenever the localized
Jacobian is nondegenerate.  In the resulting variables \((V,P)\),
\[
  \partial_\tau V+(V\cdot\nabla)V+\nabla P
  =\nu\Delta V+a(\tau)(V+y\cdot\nabla V)+b(\tau)\cdot\nabla V,
  \qquad \nabla\cdot V=0,
\]
and the pressure is determined modulo functions of time by the compact-ball
pressure decomposition.  Failure of the localized nondegeneracy condition is a
multibubble or gauge-degenerate alternative.
\end{theorem}

\begin{proof}
The localized scale and centering functionals have an invertible derivative on
the selected core.  The quantitative inverse function theorem gives the local
scale-translation choice on each compact selected window.  Differentiating the
rescaled variables gives the displayed equation, and the pressure transforms
with the usual freedom to add functions of time.  If the derivative is not
invertible, a unique core has not been selected, which is precisely the local
gauge-degenerate alternative.
\end{proof}

\begin{theorem}[Local multibubble and cascade exclusion]\label{thm:multibubble}
Assume the local Type I tangent criterion and the local multibubble/cascade
decoupling and scale-rigidity theorems.  Then no local multibubble, cascade, or
gauge-degenerate Type II configuration occurs.
\end{theorem}

\begin{proof}
If compact-cylinder bounds fail, or if active CKN mass splits, another active
bounded-cylinder concentration core is selected.  Comparable cores are combined
or remain in the multibubble case.  Strict same-point cascades are reduced to
scale-rigidity or nonlinear no-splitting.  Separated physical cores decouple in
active local frames.  The local decoupling and scale-rigidity theorems then
reduce every such configuration either to the single-core criterion or to the
Type I tangent criterion, both of which exclude it.
\end{proof}

\begin{theorem}[Compact-cylinder Caccioppoli regularity]\label{thm:rough-core}
Let \((V,P,a,b)\) solve the renormalized Navier--Stokes equation on compact
cylinders, with local pressure reconstruction modulo functions of time, bounded
critical local norm on the cylinders under consideration, and bounded or
window-integrable modulation coefficients.  Then the local Caccioppoli estimate
gives uniform local windowed \(H^1\) control.  Consequently the rough-core
alternative, meaning failure of this local windowed \(H^1\) bound in a bounded
renormalized core, cannot occur under these hypotheses.
\end{theorem}

\begin{proof}
The local energy inequality is pulled back through the absolutely continuous
scale-translation chart.  A compactly supported cutoff gives the local
Caccioppoli estimate.  The pressure term is estimated modulo constants by the
compact-ball pressure decomposition and Calderon--Zygmund/harmonic bounds.
The local \(L^2\) and \(L^3\) terms are controlled on bounded cylinders, and
the modulation terms are controlled by the stated coefficient hypotheses.  This
yields the windowed \(H^1\) bound.
\end{proof}

\begin{theorem}[Scale-collapse exclusion alternatives]\label{thm:scale-collapse}
Let \((V,P,a,b,\lambda)\) be a renormalized Type II rescaling satisfying
\[
  \partial_\tau\log\lambda=a,
  \qquad \lambda(\tau)\to0,
\]
and suppose a localized \(L^2\) core has a positive floor.  Then
\[
  \int_{\tau_0}^{\infty}a_-(\tau)M_R(\tau)\,d\tau=\infty.
\]
If the finite-cost class is controlled by this scale-drift cost, the
finite-cost scale-collapse rescaling is excluded.  Alternatively, on thick
scale-collapse windows with compactness, modulation convergence, and a
stationary omega-limit, the rescaling produces a nonzero stationary or
generalized self-similar profile.  In parameter ranges where the corresponding
Liouville theorem excludes such profiles, the scale-collapse rescaling is
excluded.
\end{theorem}

\begin{proof}
The identity \(\partial_\tau\log\lambda=a\) implies
\[
  \int_{\tau_1}^{\tau_2}a(\tau)\,d\tau
  =\log\lambda(\tau_2)-\log\lambda(\tau_1).
\]
If \(\lambda(\tau)\to0\), the right-hand side tends to \(-\infty\).  Hence
\(\int a_- =\infty\), and the positive localized \(L^2\) floor gives divergence
of \(\int a_-M_R\).  The cost conclusion follows from the finite-cost control
hypothesis.  The structural alternative follows by translating thick windows,
using compactness and modulation convergence to pass to an autonomous reduced
limit, and then applying the stationary omega-limit and Liouville hypotheses.
\end{proof}

\section{Assembly theorem}

\begin{theorem}[Local Type II exclusion theorem]\label{thm:assembly}
Let \((u,p)\) be a suitable weak Navier--Stokes solution near \((x_0,T)\).
Assume the local Type I tangent criterion, the local multibubble/cascade
decoupling and scale-rigidity theorems, the compact-cylinder Caccioppoli
hypotheses, and the scale-collapse cost or Liouville hypotheses stated in
\cref{thm:scale-collapse}.  Then \((x_0,T)\) is not a Type II singular point.
\end{theorem}

\begin{proof}
Suppose, for contradiction, that \((x_0,T)\) is a Type II singular point.  By
\cref{thm:positive-concentration}, there is a sequence of parabolic rescalings
with positive local CKN concentration.  Since the Type I tangent case is removed
by \cref{ass:typeI}, the sequence belongs to the local Type II regime.

Apply the compact-cylinder dichotomy, \cref{thm:compactness-dichotomy}.  If the
compact-cylinder suitable package fails, then the sequence enters the local
multibubble or cascade alternative, which is excluded by \cref{thm:multibubble}.
Thus we may assume compact-cylinder suitable bounds.

If the localized scale-translation conditions are degenerate, the sequence is
again in the multibubble or gauge-degenerate alternative and is excluded by
\cref{thm:multibubble}.  Otherwise \cref{thm:representation} gives the
renormalized scale-translation variables and the local pressure representation
on the selected single core.

For this renormalized single-core sequence, consider the possible failures of the
selected-window compact argument.  If the localized Type II cost is finite, then
\cref{thm:single-core} excludes the sequence.  If local windowed \(H^1\) control
fails in a bounded renormalized core, \cref{thm:rough-core} excludes that
alternative under the compact-cylinder Caccioppoli hypotheses.  If the
renormalized scale collapses while a localized core remains nonzero,
\cref{thm:scale-collapse} excludes the scale-collapse alternative under the
stated finite-cost or Liouville hypotheses.

These alternatives exhaust the positive local Type II concentration cases:
compact single-core finite cost, multibubble/cascade or gauge degeneracy,
rough-core loss of local windowed \(H^1\), and scale collapse.  Each case has
been excluded, contradicting the assumed Type II singularity.
\end{proof}

\begin{corollary}[Conditional exclusion after Type I]
Assume the local Type I tangent criterion and all analytic hypotheses listed in
\cref{thm:assembly}.  Then no singular point of a suitable weak solution can
belong to the local Type II regime.
\end{corollary}

\begin{proof}
This is \cref{thm:assembly} applied at each putative singular point after the
positive local concentration reduction.
\end{proof}

\section{Remaining analytic inputs}

The assembled theorem leaves no hidden alternatives.  Any failure of the
conclusion must occur through failure of one of the stated PDE inputs:

\begin{enumerate}[label=\textup{(\roman*)}]
\item the local Type I tangent criterion;
\item compact-cylinder suitable compactness or the corresponding local
multibubble/cascade exclusion;
\item localized scale-translation nondegeneracy on a selected core;
\item compact-cylinder Caccioppoli control and the pressure decomposition used
in the rough-core theorem;
\item finite-cost control by the localized scale-drift cost, or the
parameter-specific Liouville theorem for the stationary or generalized
self-similar scale-collapse profile.
\end{enumerate}

Thus the assembly theorem is a conditional Type II exclusion theorem in purely
PDE terms: every local Type II concentration is either one of the listed
alternatives or is excluded by the corresponding theorem.

\bibliographystyle{alpha}
\bibliography{paper6_pde_typeII_exclusion_assembly}

\end{document}
